\section{Time-Frequency Analysis}
\label{Chap:Time-Frequency Analysis}
\begin{frame}[t]{Chapter \autoref{Chap:Time-Frequency Analysis}}
\tableofcontents[sections= \value{section}]
\end{frame}


\subsection{Introduction}
\begin{frame}[t]{Introduction}

\begin{itemize}

\item Time-Frequency analysis is used to analyse non-stationary signals: signals which varies in time

\item Common Tool to this used is called Short-time-Fourier-Transform (STFT)

\end{itemize}

\end{frame}


\subsection{Short Time Fourier Transform}
\begin{frame}[t]{STFT Concepts}

\begin{itemize}

\item Instead of sliding the window on the data, we will do the opposite:

	\begin{itemize}
	\item We fix the window of size $L$
	
	\item We choose an index $n$ which represent the shift in the data $x[m]$: so $n$ represents the start of the segment we are looking on 
	
	\item We make the product between the window and the data
	
	\item Then take the DTFT of this product
	
	
	\item Then we repeat this process for every value of $n \, \rightarrow x[n] \, , \, x[n \, + \, 1] \, , \, \cdots \, x[n \, + \, L - \, 1]$
	\end{itemize}

\item Common Tool to this used is called Short-time-Fourier-Transform (STFT)

\end{itemize}


\end{frame}


\subsection{Temporal and Frequency Resolution Trade-off}
\begin{frame}[t,allowframebreaks]{Temporal and Frequency Resolution Trade-off}

\begin{itemize}

\item From window analysis, we can recall that the higher the window length $\rightarrow$ more frequency details we can capture $\rightarrow$ higher frequency resolution

	\begin{itemize}
	\item However when we do this in STFT, we loose temporal resolution $\leftrightarrow$ we can't observe details in time domain in the time-frequency plane
	\end{itemize}

\item If we decrease the window size, the resolution in frequency content decrease

	\begin{itemize}
	\item For a window size $L \, = \, 128$ Along the vertical axis (frequency axis), we can't capture that we have 2 frequencies at 250 Hz and 300 Hz
	
	\item However for this same window size, we get a better temporal resolution $\leftrightarrow$ we can indentify the moment of transition from 300 Hz $\rightarrow$ 250 Hz
	\end{itemize}


\end{itemize}


\end{frame}
